\par\noindent\textbf{(1)} 实验室Bloch方程为 Ṁ=γ M×B−(M\_xx̂+M\_yŷ)/T₂−(M\_z−M₀)ẑ/T₁。取绕z角速度ω的旋转系，有有效场 B\_eff=B₁x̂+(B₀−ω/γ)ẑ，并按γ和坐标约定确定转矩符号。定义 Δ=γB₀−ω、Ω₁=γB₁。
\par\noindent\textbf{(2)} 在弱驱动、M\_z≈M₀的线性稳态，横向两个方程可写 0=−M\_x/T₂−ΔM\_y，0=ΔM\_x−M\_y/T₂−Ω₁M₀。解得 M\_y=−Ω₁M₀ T₂/[1+Δ²T₂²]、M\_x=Ω₁M₀ΔT₂²/[1+Δ²T₂²]（符号取决于驱动旋向定义，耗散功率为正）。
\par\noindent\textbf{(3)} 射频场平均吸收功率密度的正值与 Ω₁|M\_y| 成正比，因此频率线型 A(Δ)∝T₂/[1+Δ²T₂²]。在 |Δ|=1/T₂ 时为峰值一半，角频率FWHM为2/T₂；T₂控制横向相位记忆和线宽，而非T₁。
\par\noindent\textbf{(4)} 当射频不再弱，可从纵向稳态得到 M\_z=M₀/[1+Ω₁²T₁T₂/(1+Δ²T₂²)]，饱和因子s=Ω₁²T₁T₂。s≪1才允许M\_z≈M₀；B₁→0恢复平衡磁化且吸收按B₁²消失。静态不均匀展宽需另卷积，不能单用T₂解释全部谱宽。
